\chapter{Sums of Squares}\label{ch:squares}

\begin{center}\itshape
$2026=1^2+45^2$, $2027$ is not a sum of two squares, and every non-negative integer is a sum of four.\\
Which integers are sums of two squares is decided by a single congruence on their prime factors.
\end{center}


The number $2026$ is a sum of two squares: $2026=1^2+45^2$. The next integer is different. Every integer square is $0$ or $1$ modulo $4$, so a sum of two squares is $0$, $1$ or $2$ modulo $4$. Since $2027\equiv3\pmod4$, it cannot be a sum of two integer squares. A short congruence has decided a question that a long search would not settle.

But passing this test is not enough. A number can be $1$ modulo $4$ and still fail to be a sum of two squares. The full theorem looks at its prime factorisation: every prime congruent to $3$ modulo $4$ must occur with an even exponent. The condition is both necessary and sufficient. Proving each direction is a different job.

A useful number to keep in mind is $65$. It has two representations, $65=1^2+8^2=4^2+7^2$. The identities are easy to check. Explaining the number of representations requires multiplication in the Gaussian integers, which we will develop in Chapter~\ref{ch:gaussian}. First we earn the two-square criterion using tools already available.

Four squares tell a different story. No congruence prevents an integer from being a sum of four squares: Lagrange's theorem says every nonnegative integer has such a representation. The proof must produce existence for all inputs, not merely classify a subset. We will separate the two problems and keep the integer zero in their domains.


\portrait{fermat}{Pierre de Fermat (1601--1665). Seventeenth-century portrait, artist unknown. Public domain, via Wikimedia Commons.}

\section{Two squares}

% t1-pass2 motive 8fef34db
Test a few numbers before reaching for the theorem. We have $5=1^2+2^2$ and $13=2^2+3^2$, but $7$ is impossible by reducing squares modulo $4$. That first obstruction is quick and useful, yet it does not classify every integer. The complete two-square criterion looks at the exponents of prime factors congruent to $3$ modulo $4$. We will prove both directions: every representation obeys the condition, and every integer obeying it has a representation.

The \emph{Brahmagupta--Fibonacci identity} shows that sums of two squares are closed under products:
\[
(a^2+b^2)(c^2+d^2)=(ac-bd)^2+(ad+bc)^2,
\]
which is $|z|^2|w|^2=|zw|^2$ for $z=a+ib$, $w=c+id$.

\begin{lemma}[Thue]\label{lem:thue}
Let $p$ be prime and $p\nmid a$. There are integers $x,y$ with $0<|x|,|y|<\sqrt p$ and $ax\equiv y\pmod p$.
\end{lemma}
\begin{proof}
Let $m=\lfloor\sqrt p\rfloor$.
Since a prime is not a perfect square,
$m<\sqrt p<m+1$, hence
$(m+1)^2>p$. Consider the $(m+1)^2$
pairs $(u,v)$ with $0\le u,v\le m$.
There are only $p$ residue classes modulo $p$,
so two distinct pairs $(u,v)$ and $(u',v')$
give the same residue of $u-av$:
\[
u-av\equiv u'-av'\pmod p.
\]
Subtract and set $x=v-v'$, $y=u-u'$.
Then
\[
ax\equiv y\pmod p,\qquad |x|,|y|\le m<\sqrt p.
\]
We still have to show that neither difference
is zero. If $x=0$, the congruence says
$p\mid y$. Since $|y|<\sqrt p\le p$,
we get $y=0$, contradicting the fact
that our pairs were distinct.
If $y=0$, then $p\mid ax$; the hypothesis
$p\nmid a$ implies $p\mid x$, and the
same size bound forces $x=0$, again a
contradiction. Thus both $x$ and $y$
are nonzero, as required.
\end{proof}

\begin{think}
Which primes below $50$ are sums of two squares? Write them down, then reduce them modulo $4$. Why can a prime $\equiv3\pmod4$ never work? The other direction is the real theorem: what would you need to build $x$ and $y$ from $p$?
\end{think}

\begin{theorem}[Fermat, 1640; Euler, 1749]\label{thm:twosq}
An odd prime $p$ is a sum of two squares if and only if $p\equiv1\pmod4$.
\end{theorem}
\begin{proof}
Squares are $0$ or $1$ modulo $4$, so a sum of two squares is never $\equiv3\pmod4$. Conversely, if $p\equiv1\pmod4$, $-1\equiv a^2$ for some $a$ (Chapter~\ref{ch:reciprocity}). Take $x,y$ from Thue's lemma: $y^2\equiv a^2x^2\equiv-x^2$, so $p\mid x^2+y^2$, and $0<x^2+y^2<2p$. Hence $x^2+y^2=p$.
\end{proof}

% t1-pass2 why
\paragraph{Why this works.}
For a prime $p\equiv1\pmod4$, reciprocity gives a square root $a$ of $-1$ modulo $p$. Thue's lemma finds small nonzero integers $x,y$ with $y\equiv ax\pmod p$. Thus $x^2+y^2\equiv(a^2+1)x^2\equiv0\pmod p$. Since both squares are smaller than $p$, their positive sum is less than $2p$. It is a positive multiple of $p$, so it equals $p$. The reverse direction is the easy check of square residues modulo $4$.

\begin{theorem}\label{thm:twosqn}
An integer $n\ge1$ is a sum of two squares if and only if every prime $q\equiv3\pmod4$ appears in $n$ with an even exponent.
\end{theorem}
\begin{proof}
If the condition holds, $n=m^2\cdot2^a\cdot\prod p_i$ with $p_i\equiv1\pmod4$; each factor is a sum of two squares ($m^2=m^2+0^2$, $2=1^2+1^2$), and so is the product. Conversely, let $n=x^2+y^2$ and $q\equiv3\pmod4$ divide $n$. If $q\nmid x$, then $\left(yx^{-1}\right)^2\equiv-1\pmod q$, impossible. So $q\mid x$, $q\mid y$, $q^2\mid n$, and $\frac n{q^2}=\left(\frac xq\right)^2+\left(\frac yq\right)^2$. Induction on $v_q(n)$ shows $v_q(n)$ is even.
\end{proof}

% t1-pass2 why
\paragraph{Why this works.}
The two-square identity tells us how to multiply known representations. Squares themselves are represented trivially, and primes congruent to $1$ modulo $4$ were handled above. What about a prime $q\equiv3\pmod4$ dividing $x^2+y^2$? If it did not divide both coordinates, $-1$ would be a square modulo $q$, which is impossible. So it divides both, and $q^2$ can be removed from the represented integer. Repeating this shows why those primes must have even exponents.

\subsection*{From a theorem to a working recipe}
The two-square criterion is an existence test. Given an integer, factor it and look at primes congruent to $3$ modulo $4$. If any such prime has odd exponent, stop: no two-square representation exists. If all their exponents are even, the theorem promises a representation. Finding the actual two integers takes another step: represent the primes congruent to $1$ modulo $4$ and multiply those representations with the Brahmagupta--Fibonacci identity. A yes-or-no theorem and a construction are related, but they are not the same output.

For instance,
\[
45=3^2\cdot5=3^2(1^2+2^2)=3^2+6^2.
\]
The prime $3$ causes no obstruction here because its exponent is even. By contrast $75=3\cdot5^2$ is impossible: adding a large square factor does not cure an odd exponent of $3$. The number $65=5\cdot13$ illustrates another feature. Multiplying $5=1^2+2^2$ and $13=2^2+3^2$ in the two possible sign choices gives
\[
\begin{aligned}
65&=(1\cdot2-2\cdot3)^2+(1\cdot3+2\cdot2)^2=4^2+7^2,\\
65&=(1\cdot2+2\cdot3)^2+(1\cdot3-2\cdot2)^2=8^2+1^2.
\end{aligned}
\]
Fermat's uniqueness for a prime, understood up to signs and order, therefore does not extend to composite numbers. Primes are choosy; products have a social life.

\begin{figure}[ht]
\centering
\fig{ch07-sixtyfive}
\caption{The first quadrant of $x^2+y^2=65$. Its integer points are $(1,8)$, $(4,7)$, $(7,4)$ and $(8,1)$; reflecting them across the axes gives $16$ representations when signs and order count. None lies on an axis, so each reflection creates a new point.}
\label{fig:sixtyfive}
\end{figure}

\begin{reminder}[A prime dividing a sum need not divide either term]
From $5\mid1^2+2^2$ we cannot infer $5\mid1$ or $5\mid2$. What makes a prime $q\equiv3\pmod4$ special is the absence of a square root of $-1$. The proof uses that extra fact to force divisibility of \emph{both} coordinates, then removes $q^2$. This is a local obstruction, not the general rule for divisibility of sums.
\end{reminder}

Counting representations is subtler. Let $r_2(n)$ be the number of $(x,y)\in\Z^2$ with $x^2+y^2=n$ (signs and order count). Then
\[
\begin{gathered}r_2(n)=4\sum_{d\mid n}\chi_4(d),\qq\chi_4(d)=\begin{cases}0&d\text{ even},\\1&d\equiv1\pmod4,\\-1&d\equiv3\pmod4,\end{cases}\end{gathered}
\]
a formula of Jacobi. \emph{We do not use it here; it is proved in Chapter~\ref{ch:gaussian} with Gaussian integers.} For instance $r_2(25)=4(1+1+1)=12$: $(\pm5,0),(0,\pm5),(\pm3,\pm4),(\pm4,\pm3)$.

% t1-pass2 examples 8fef34db
\begin{examplex}[the integers up to $20$]
A positive integer is a sum of two squares precisely when every prime $3$ modulo $4$ has even exponent. Among $1,\ldots,20$ the successes are
\[
1,2,4,5,8,9,10,13,16,17,18,20.
\]
For instance $18=3^2+3^2$ and the prime $3$ has even exponent, while $12=4\cdot 3$ is rejected. Figure~\ref{fig:t1-representable} records the same list through $30$.
\end{examplex}

\begin{figure}[ht]
\centering
\fig{ch07-representable}
\caption{Integers from $1$ to $30$. Dark grey cells are sums of two squares; rust cells are not. The failures include every number congruent to $3$ modulo $4$, and also $12=4\cdot 3$, where the prime $3$ has odd exponent.}
\label{fig:t1-representable}
\end{figure}

\begin{examplex}[building $65$ again, from the identity]
One has $5=1^2+2^2$ and $13=2^2+3^2$. The two sign choices give $1^2+8^2$ and $4^2+7^2$, both equal to $65$. A prime $1$ modulo $4$ has only one representation up to signs and order; a product of two such primes is the first place a second pair appears.
\end{examplex}

\section{Four squares}

% t1-pass2 motive 85e01d27
The number $7$ defeats the two-square question, but $7=2^2+1^2+1^2+1^2$. In fact every nonnegative integer is a sum of four integer squares. That is Lagrange's theorem. To prove it for every integer, a few successful decompositions are not enough. We need an identity that preserves sums of four squares under multiplication, then a descent that gives representations of primes. Only after those ingredients are proved may we assemble arbitrary integers by factorisation.

\begin{think}
Find the integers $n\le30$ that are not sums of three squares. What do they have in common modulo $8$? Do any of them need more than four squares?
\end{think}

\begin{theorem}[Lagrange, 1770]\label{thm:foursq}
Every positive integer is a sum of four squares.
\end{theorem}
\begin{proof}
\emph{Euler's identity.} One checks by expansion (or with quaternion norms, Volume~3) that
\[
\left(\sum_{i=1}^4a_i^2\right)\left(\sum_{i=1}^4b_i^2\right)=z_1^2+z_2^2+z_3^2+z_4^2,
\]
with $z_1=a_1b_1+a_2b_2+a_3b_3+a_4b_4$, $z_2=a_1b_2-a_2b_1+a_3b_4-a_4b_3$, $z_3=a_1b_3-a_3b_1+a_4b_2-a_2b_4$, $z_4=a_1b_4-a_4b_1+a_2b_3-a_3b_2$. So it suffices to treat primes, and $2=1+1+0+0$.

\emph{A first multiple.} Let $p$ be an odd prime. The $\frac{p+1}2$ numbers $x^2$, $0\le x\le\frac{p-1}2$, are distinct modulo $p$, and so are the $\frac{p+1}2$ numbers $-1-y^2$. Together they are $p+1$ numbers, so some $x^2\equiv-1-y^2$: $x^2+y^2+1+0^2=mp$ with $0<mp\le2\left(\frac{p-1}2\right)^2+1<p^2$, so $1\le m<p$.

\emph{Descent.} Let $m$ be the least positive integer with $mp=x_1^2+x_2^2+x_3^2+x_4^2$; we know $m<p$. If $m$ is even, an even number of the $x_i$ are odd; pair them so that $x_1\equiv x_2$, $x_3\equiv x_4\pmod2$, and $\frac m2p=\left(\frac{x_1+x_2}2\right)^2+\left(\frac{x_1-x_2}2\right)^2+\left(\frac{x_3+x_4}2\right)^2+\left(\frac{x_3-x_4}2\right)^2$ contradicts minimality. So $m$ is odd. Suppose $m>1$. Choose $y_i\equiv x_i\pmod m$ with $|y_i|\le\frac{m-1}2$. Then $\sum y_i^2\equiv\sum x_i^2\equiv0\pmod m$, say $\sum y_i^2=mr$, with $r\le\frac{4}{m}\left(\frac{m-1}2\right)^2<m$. And $r>0$: otherwise all $x_i\equiv0\pmod m$, so $m^2\mid mp$, $m\mid p$, impossible for $1<m<p$.
Now apply Euler's identity to $(mp)(mr)$ with $a=x$, $b=y$. Since $y_i\equiv x_i\pmod m$: $z_1\equiv\sum x_i^2\equiv0$, and each of $z_2,z_3,z_4$ is $\equiv0$ because its terms cancel in pairs modulo $m$ (e.g.\ $x_1y_2-x_2y_1\equiv x_1x_2-x_2x_1$). So $rp=\sum\left(\frac{z_i}m\right)^2$ with $0<r<m$, contradicting minimality. Hence $m=1$.
\end{proof}

% t1-pass2 why
\paragraph{Why this works.}
The four-square identity makes representations stable under multiplication, so the hard work reduces to primes. For an odd prime $p$, a pigeonhole argument first produces a representation of some positive multiple $mp$ with $m<p$. The descent uses the same four-square identity and a congruence modulo $m$ to replace that multiplier by a smaller positive one. A strictly decreasing sequence of positive multipliers must end. It can end only at $m=1$, which is exactly the representation we wanted.

\begin{remark}
Three squares are harder: \emph{Legendre and Gauss} proved that $n$ is a sum of three squares if and only if $n$ is not of the form $4^a(8b+7)$. The ``only if'' direction is Exercise~\ref{ex:three}; the converse needs either Gauss' theory of ternary quadratic forms or Dirichlet's theorem on primes in progressions, and is proved in Volume~13 (with the Hasse--Minkowski theorem). We only use it in Exercise~\ref{ex:eureka}, where it is explicitly admitted.
\end{remark}

\subsection*{Watching the descent actually descend}
The four-square proof has three jobs: produce \emph{some} multiple of $p$, make a smaller multiplier, and stop at $1$. The first pigeonhole argument supplies a starting point; the identity supplies a machine for moving down. Minimality turns the machine into a proof without making us run it forever.

Here is one run with $p=7$. Start, deliberately inefficiently, from
\[
3\cdot7=4^2+2^2+1^2+0^2.
\]
The nearest residues modulo $m=3$ are $(1,-1,1,0)$, whose squares sum to $3=mr$ with $r=1$. Euler's identity gives $(z_1,z_2,z_3,z_4)=(3,-6,3,3)$. Dividing the four coordinates by $3$ yields
\[
7=1^2+(-2)^2+1^2+1^2.
\]
The multiplier really fell from $3$ to $1$. It is essential that we divide the \emph{coordinates} by $m$ and the squared norm by $m^2$: $(mp)(mr)/m^2=pr$. Losing one of these factors is an impressively efficient way to prove something false.

Why do we remove even multipliers first? With odd $m$, nearest residues satisfy $|y_i|\le(m-1)/2$, so the new multiplier is \emph{strictly} smaller. Using only $|y_i|\le m/2$ would give $r\le m$, leaving open the possibility that the machine merely stands still.

\begin{think}
Can $7$ be a sum of three squares? Squares modulo $8$ are $0,1,4$; check every possible sum of three such residues. Yet $7=2^2+1^2+1^2+1^2$. This is a small certificate that ``four'' in Lagrange's theorem cannot be replaced by ``three''.
\end{think}

% t1-pass2 examples 85e01d27
\begin{examplex}[$7$ needs four squares]
Modulo $8$, squares are $0,1,4$, so sums of three squares miss $7$. Thus $7$ is not a sum of three squares. It is a sum of four: $4+1+1+1$. The identity is not needed for so small a number; it is needed to know that every larger integer inherits a representation from its prime factors.
\end{examplex}

\begin{examplex}[a composite through its prime factors]
One has $3=1^2+1^2+1^2+0^2$ and $5=2^2+1^2+0+0$. Their product $15$ must be a sum of four squares by Euler's identity; an explicit representation is $15=3^2+2^2+1^2+1^2$. The identity guarantees that some representation exists before one goes looking for the squares.
\end{examplex}

\section{Lattice points in a disc}

% t1-pass2 motive 022b77aa
A solution of $a^2+b^2=n$ is a point on a circle of radius $\sqrt n$ in the integer lattice. Count solutions for every $n$ up to $R^2$, and you are counting \emph{all} lattice points inside a disc of radius $R$. Draw a unit square around each point: the union resembles a disc, with the disagreement concentrated near its boundary. This gives the main term $\pi R^2$ and a first error bound. A sharper error needs a sharper idea.

\begin{think}
How many points of $\Z^2$ lie in the disc of radius $10$ centred at $0$? Guess first, using only the area of the disc. (The answer is $317$, against $100\pi\approx314.2$.) Why should the area be a good guess, and how large could the error be?
\end{think}

\begin{theorem}[Gauss]\label{thm:circle}
Let $N(R)=\#\{(x,y)\in\Z^2:x^2+y^2\le R^2\}$. Then $\left|N(R)-\pi R^2\right|\le2\sqrt2\,\pi R+2\pi$ for $R\ge1$. In particular $N(R)\sim\pi R^2$.
\end{theorem}
\begin{proof}
Attach to each counted point $(x,y)$ the unit square $[x,x+1[\times[y,y+1[$. These squares are disjoint, have total area $N(R)$, and lie in the disc of radius $R+\sqrt2$ (a point of a square is within $\sqrt2$ of its corner). They cover the disc of radius $R-\sqrt2$: a point $P$ of that disc lies in the square of the lattice point $(\lfloor P_x\rfloor,\lfloor P_y\rfloor)$, which is within $\sqrt2$ of $P$, hence in the disc of radius $R$. Thus $\pi\max(0,R-\sqrt2)^2\le N(R)\le\pi(R+\sqrt2)^2$. For $R\ge\sqrt2$, expanding gives the claimed bound. For $1\le R<\sqrt2$, the desired lower estimate $N(R)\ge\pi R^2-2\sqrt2\,\pi R-2\pi$ is automatic because its right-hand side is negative; the upper estimate follows from the same outer disc. The maximum matters: a disc cannot have negative radius.
\end{proof}

% t1-pass2 why
\paragraph{Why this works.}
Give each integer lattice point a unit square. Those squares have total area equal to the number of points, because their interiors do not overlap. The union contains a slightly smaller disc and lies inside a slightly larger disc. The areas of those two discs differ by a boundary strip of fixed width, proportional to the radius $R$. That sandwich gives the leading term $\pi R^2$ and an error of order $R$. The picture supplies a bound, not the exact number of points.

\begin{figure}[ht]
\centering
\fig{ch07-disc}
\caption{The $N(6)=113$ lattice points in the disc of radius $6$, against $36\pi\approx113.1$. Each point $(x,y)$ owns the unit square $[x,x+1[\times[y,y+1[$. The squares cover the disc of radius $R-\sqrt2$ and stay inside the disc of radius $R+\sqrt2$ (dashed), which squeezes $N(R)$ between the two areas.}
\label{fig:disc}
\end{figure}

Since $N(R)=1+\sum_{1\le n\le R^2}r_2(n)$, where the extra $1$ counts the origin, the \emph{average} of $r_2(n)$ tends to $\pi$: explicitly, $X^{-1}\sum_{1\le n\le X}r_2(n)\to\pi$ as $X\to\infty$. The true size of the error is the \emph{Gauss circle problem}: it is believed to be $O(R^{1/2+\varepsilon})$, and the cited work of Bourgain and Watt (2017) gives a bound $O(R^{517/824+\varepsilon})$, about $R^{0.6275+\varepsilon}$. This is a dated benchmark, not a claim that it remains the latest record.

\subsection*{What the area argument does, and does not, say}
Draw unit squares around the lattice points inside the disc. Away from the boundary, the squares behave exactly as the area count suggests: whole squares are included or excluded. Only squares that meet the circumference can change the answer. A strip of fixed width around a circle of radius $R$ has area proportional to $R$, so the discrepancy is at most of that order. This gives a useful first bound without proving that every boundary square contributes the same way.

This does not mean that $N(R)-\pi R^2$ always has the same sign, or that it is small for every tiny radius. The staircase $N(R)$ jumps only when $R^2$ passes a sum of two squares, while $\pi R^2$ changes smoothly. The theorem controls the discrepancy without pretending the staircase is a curve. Dividing its bound by $\pi R^2$ gives a relative error tending to zero, which is precisely the assertion $N(R)\sim\pi R^2$.

% t1-pass2 examples 022b77aa
\begin{examplex}[the disc of radius $2$]
The lattice points with $x^2+y^2\le 4$ are the origin, the four points at distance $1$, and the four points at distance $\sqrt 2$. So $N(2)=9$. The area is $4\pi\approx 12.57$. The error bound of the theorem is $2\sqrt 2\,\pi\cdot 2+2\pi\approx 24$, which is true and deliberately crude: the argument is shaped to give the order of the error, not the sharp constant.
\end{examplex}

\begin{examplex}[why the squares stick out]
The unit square to the north-east of $(2,0)$ reaches the point $(3,1)$, at distance $\sqrt{10}>3$ if one were looking at radius $2$. The proof therefore places the squares inside a disc of radius $R+\sqrt 2$, not of radius $R$. That is the whole of the linear error.
\end{examplex}

\begin{beauty}
Where this leads. The two-square theorem becomes unique factorisation in $\Z[i]$ in Chapter~\ref{ch:gaussian}, and the class number of $\Q(i)$ is $1$. The three-square theorem and the theta series that packages $r_2(n)$ are taken up in Volumes~13 and~6. The error in the lattice-point count is a model of the exponential-sum problems in Volume~14.
\end{beauty}

\subsection*{Proof workshop: an obstruction is not a classification}

Start with $n=7$. If $7=a^2+b^2$ for integers $a,b$,
then reduction modulo $4$ would express $3$ as a sum
of two terms from $\{0,1\}$. That is impossible. The same
argument rules out \emph{every} integer congruent to
$3$ modulo $4$.

What does it say about $21$? Again $21\equiv1\pmod4$,
so the argument is silent. But $21=3\cdot7$ has two
primes congruent to $3$ modulo $4$, each with odd
exponent. The two-square theorem rules it out.
A necessary condition such as a modular test is
useful even when it is not sufficient. To classify
all solutions, the proof must also establish the
converse of its final criterion.


\section{Exercises}

\begin{exercise}[$\star$]\label{ex:three}
For $a,b\in\N$, show that the integer $4^a(8b+7)$ is not a sum of three squares. \emph{Look at squares modulo $8$.}
\end{exercise}

\begin{exercise}[$\star$]\label{ex:2026}
Write $1013$ (a prime) and then $2026$ as sums of two squares. Explain why $2027$ is not a sum of two squares, and find all $n\in\{2021,\dots,2030\}$ that are.
\end{exercise}

\begin{exercise}[$\star\star$]\label{ex:primitive2}
Show that if $n=x^2+y^2$ with $\gcd(x,y)=1$, then $4\nmid n$ and $n$ has no prime factor $\equiv3\pmod4$.
\end{exercise}

\begin{exercise}[$\star\star$]\label{ex:x2y2}
Show that a prime $p\equiv1$ or $3\pmod8$ is of the form $x^2+2y^2$. \emph{Use Thue's lemma with $a^2\equiv-2$ (Chapter~\ref{ch:reciprocity}), then deal with the case $x^2+2y^2=2p$.}
\end{exercise}

\begin{exercise}[$\star\star$]\label{ex:x3y2}
Show that a prime $p\equiv1\pmod3$ is of the form $x^2+3y^2$.
\end{exercise}

\begin{exercise}[$\star\star$]\label{ex:eureka}
\emph{Gauss' ``E\upshape{}YPHKA'' (1796).} Admitting the three-square theorem, show that every $8n+3$ is a sum of three odd squares, and deduce that every $n\in\N$ is a sum of three triangular numbers $\frac{k(k+1)}2$.
\end{exercise}

\begin{exercise}[$\star\star\star$]\label{ex:zagier}
\emph{Zagier's one-sentence proof (1990).} Let $p\equiv1\pmod4$ be prime and $S=\{(x,y,z)\in(\N^*)^3:x^2+4yz=p\}$.
\begin{enumerate}[label=(\alph*)]
\item Show that $S$ is finite and that $x\ne y-z$ and $x\ne2y$ for $(x,y,z)\in S$.
\item Show that the following map is an involution of $S$ (it sends $S$ to $S$ and $f\circ f=\operatorname{id}$):
\[
\begin{aligned}&x<y-z:\eqbr\quad f=(x+2z,\ z,\ y-x-z),\\&y-z<x<2y:\eqbr\quad f=(2y-x,\ y,\ x-y+z),\\&x>2y:\eqbr\quad f=(x-2y,\ x-y+z,\ y).\end{aligned}
\]
\item Show that $f$ has exactly one fixed point. Deduce that $|S|$ is odd.
\item Deduce that the involution $(x,y,z)\mapsto(x,z,y)$ has a fixed point, and that $p$ is a sum of two squares.
\end{enumerate}
\end{exercise}

\begin{exercise}[$\star$]\label{ex:t1-which20}
Which of $15,18,45,50$ are sums of two squares? For each answer give either a representation or the prime that forbids it.
\end{exercise}

\begin{beauty}
Two squares, three squares, four squares: the answers are, respectively, a condition on primes $\equiv3\pmod4$, a condition modulo $8$, and ``always''. The two-square theorem becomes transparent in the Gaussian integers (Chapter~\ref{ch:gaussian}), where $p=x^2+y^2$ means $p=(x+iy)(x-iy)$; the four-square theorem in the quaternions (Volume~3). The study of which primes are $x^2+ny^2$ grows into class field theory, and Zagier's proof is a reminder that sometimes a clever involution is all one needs.
\end{beauty}

\section{Hints}

\begin{hint}{ex:three}
List the residues of squares modulo $8$: only $0,1,4$ occur. Can three of them sum to $7$ modulo $8$? To handle the factor $4^a$, first show that a sum of three squares divisible by $4$ forces every variable to be even. Then divide by $4$ and repeat.
\end{hint}
\begin{hint}{ex:2026}
Factor each number first. A prime congruent to $3$ modulo $4$ with odd exponent is an obstruction; otherwise the theorem promises a representation. To find actual coordinates, represent the primes congruent to $1$ modulo $4$ and multiply the resulting sums of two squares. Check each proposed representation by squaring.
\end{hint}
\begin{hint}{ex:primitive2}
If a prime $q\equiv3\pmod4$ divides $x^2+y^2$ and $q\nmid x$, dividing by $x^2$ produces a square root of $-1$ modulo $q$, which is impossible. So $q$ divides both coordinates, contradicting primitiveness. For the factor $4$, look directly at the possible square residues modulo $4$.
\end{hint}
\begin{hint}{ex:x2y2}
Begin with $a^2\equiv-2\pmod p$ and use Thue's lemma to obtain small nonzero $x,y$ with $p\mid y^2+2x^2$. The size bound makes this number either $p$ or $2p$. In the second case, parity lets you rearrange it as the desired form with smaller coordinates.
\end{hint}
\begin{hint}{ex:x3y2}
Here the same method begins from $a^2\equiv-3\pmod p$. Thue's lemma gives $0<y^2+3x^2<4p$, so the represented number is $p$, $2p$ or $3p$. Rule out the middle case with a congruence and use the factor of $3$ in the final case.
\end{hint}
\begin{hint}{ex:eureka}
$8T_k+1=(2k+1)^2$. If $8n+3=a^2+b^2+c^2$, look modulo $8$ to see that $a,b,c$ must all be odd.
\end{hint}
\begin{hint}{ex:zagier}
(a) $x\ne y-z$ and $x\ne2y$ because $p$ is prime and not a square. (c) A fixed point must lie in the middle case with $x=y$; then $x(x+4z)=p$. (d) An involution of a finite set has as many fixed points as the size of the set, modulo $2$. \emph{A second nudge for (b):} if $x<y-z$, the first image has $x'=x+2z$ and $y'=z$, so $x'>2y'$; the third branch must bring it back. Check branch membership before composing the formulas.
\end{hint}

\begin{hint}{ex:t1-which20}
Factor them. A prime $3$ modulo $4$ is allowed only to an even power. The prime $2$ and the primes $1$ modulo $4$ are always allowed.
\end{hint}

\section{Solutions}

\begin{solution}{ex:three}
\idea A three-square obstruction should begin with the smallest useful modulus. Every square is $0$, $1$ or $4$ modulo $8$, and an odd square is always $1$. Try to express $7$ as a sum of three such residues; it cannot be done. If an integer of the form $4^a(8b+7)$ were a sum of three squares, reduction modulo $4$ would force all three variables even. Divide their squares by $4$ and repeat until the forbidden remainder $7$ appears. That proves the obstruction, not the converse theorem.
\writeup A square is congruent to $0$, $1$ or $4$ modulo $8$. Three numbers from this set can add to any residue from $0$ through $6$, but not to $7$. Consequently $8b+7$ cannot be a sum of three integer squares.

Now suppose that $4^a(8b+7)=x^2+y^2+z^2$ for some $a\ge1$. Modulo $4$, each square is $0$ or $1$. The only way three such residues add to $0$ modulo $4$ is for all three to be $0$; hence $x,y,z$ are even. Dividing their squares by $4$ gives
\[
4^{a-1}(8b+7)=(x/2)^2+(y/2)^2+(z/2)^2.
\]
Repeat this argument $a$ times. We would obtain a three-square representation of $8b+7$, which the first paragraph ruled out. This proves the obstruction for every $a,b\in\N$.
\end{solution}

\begin{solution}{ex:2026}
\idea Use the two-square criterion as a decision procedure, then find actual coordinates. Factor the proposed integer and inspect the exponents of primes congruent to $3$ modulo $4$. If they are all even, represent the remaining factors as sums of two squares and multiply the representations with the two-square identity. At the end, square your proposed coordinates and check the original number. For $1013$, a search near $\sqrt{1013/2}$ may find the first decomposition.
\writeup First find a representation of the prime $1013$:
\[
1013=22^2+23^2.
\]
Since $2=1^2+1^2$, the Brahmagupta--Fibonacci identity gives
\[
2026=2\cdot1013=(22-23)^2+(22+23)^2=1^2+45^2.
\]
On the other hand, $2027\equiv3\pmod4$. Such an integer cannot be a sum of two squares, regardless of whether it is prime, because the possible square residues modulo $4$ are only $0$ and $1$.

To classify the full range, apply the prime-factor criterion. The numbers
\[
\begin{gathered}
2021=43\cdot47,\quad 2022=2\cdot3\cdot337,\quad
2023=7\cdot17^2,\\
2024=2^3\cdot11\cdot23,\quad
2028=2^2\cdot3\cdot13^2,\quad
2030=2\cdot5\cdot7\cdot29
\end{gathered}
\]
each contain a prime congruent to $3$ modulo $4$ to an odd power, so they are excluded. We also excluded $2027$ above. The remaining successes are
\[
2025=45^2+0^2,\qquad
2026=1^2+45^2,\qquad
2029=2^2+45^2.
\]
Each equality provides a direct check of existence.
\end{solution}

\begin{solution}{ex:primitive2}
\idea A primitive representation $n=x^2+y^2$ means $\gcd(x,y)=1$. If a prime $q\equiv3\pmod4$ divided $n$, the two-square proof forces $q$ to divide \emph{both} $x$ and $y$, which contradicts primitiveness. So every prime of that residue class is excluded. The prime $2$ requires a separate parity check. Keep the argument about prime divisibility distinct from the final classification of allowed $n$.
\writeup Suppose $n=x^2+y^2$ with $\gcd(x,y)=1$. If $4\mid n$, then modulo $4$ the squares $x^2,y^2$ are each $0$ or $1$ and must both be $0$. This makes $x,y$ even, contradicting the gcd condition. Thus $4\nmid n$.

Next let $q\equiv3\pmod4$ be a prime divisor of $n$. If $q\nmid x$, the congruence
\[
x^2+y^2\equiv0\pmod q
\]
can be divided by $x^2$, giving $(yx^{-1})^2\equiv-1\pmod q$. Quadratic reciprocity says $-1$ is not a square for such a prime. Therefore $q\mid x$. The displayed congruence then forces $q\mid y^2$, hence $q\mid y$. This again contradicts primitiveness. No such prime may divide $n$.
\end{solution}

\begin{solution}{ex:x2y2}
\idea The two-square proof starts with a square root of $-1$ modulo a prime. To represent a prime as $x^2+2y^2$, change the starting congruence to $a^2\equiv-2\pmod p$. Quadratic reciprocity tells us for which primes such a root exists. Thue's lemma then gives small integers $x,y$ with $p\mid x^2+2y^2$. Use their size bounds to determine which multiple of $p$ is possible; divisibility alone does not yet give equality.
\writeup For the stated prime $p$, quadratic reciprocity gives a residue $a$ with $a^2\equiv-2\pmod p$. Apply Thue's lemma to find nonzero integers $x,y$, each smaller in absolute value than $\sqrt p$, such that $y\equiv ax\pmod p$. Hence
\[
y^2+2x^2\equiv(a^2+2)x^2\equiv0\pmod p.
\]
This integer is positive and smaller than $3p$. It must therefore be $p$ or $2p$.

If it equals $p$, the desired representation is already found. If it equals $2p$, then $y^2$ is even, so $y=2y'$ for some integer $y'$. Divide
$4y'^2+2x^2=2p$ by $2$ to obtain
\[
p=x^2+2y'^2.
\]
Both cases yield the representation. The final parity adjustment is the step that turns a small multiple of $p$ into $p$ itself.
\end{solution}

\begin{solution}{ex:x3y2}
\idea The same pattern works with $-3$, but it has an extra case. First use reciprocity to find the primes for which $-3$ is a square. Then apply Thue's lemma to get a small positive multiple of $p$ represented by $x^2+3y^2$. Its size leaves a few possible multiples. Examine each: one may be reduced to $p$, another ruled out by parity. Only after the case split have you proved the representation theorem.
\writeup Because $p\equiv1\pmod3$, the supplementary reciprocity calculation gives a square root $a$ of $-3$ modulo $p$. Thue's lemma produces nonzero integers $x,y$ with $|x|,|y|<\sqrt p$ and $y\equiv ax\pmod p$. Therefore
\[
0<y^2+3x^2<4p,\qquad
p\mid y^2+3x^2.
\]
The left side is $p$, $2p$ or $3p$.

If it is $p$, we are done. If it is $3p$, reduction modulo $3$ shows $3\mid y$, say $y=3y'$, and division by $3$ gives $p=x^2+3y'^2$.

It cannot be $2p$. Indeed, $y^2+3x^2$ is even only if $x,y$ have the same parity. If both are even, the sum is divisible by $4$. If both are odd, their squares are $1$ modulo $4$, and the sum is again $1+3=0$ modulo $4$. But $2p$ is not divisible by $4$ because $p$ is odd. Thus only the two successful cases remain.
\end{solution}

\begin{solution}{ex:eureka}
\idea Triangular numbers are linked to odd squares by $8T_k+1=(2k+1)^2$. So ``$n$ is a sum of three triangular numbers'' is the same as ``$8n+3$ is a sum of three \emph{odd} squares''. The three-square theorem gives three squares; a look modulo $8$ forces them to be odd.
\writeup We may use the three-square theorem as granted in the question. The positive integer $8n+3$ is not of the forbidden shape $4^a(8b+7)$: it is odd and congruent to $3$ modulo $8$. Thus there are integers $a,b,c$ with
\[
8n+3=a^2+b^2+c^2.
\]
Modulo $8$, the only way three squares from $\{0,1,4\}$ can add to $3$ is $1+1+1$. Hence $a,b,c$ are all odd. Their signs do not matter, so replace them by positive absolute values and write $a=2r+1$, $b=2s+1$, $c=2t+1$ with $r,s,t\in\N$.

The identity $(2r+1)^2=8\,r(r+1)/2+1$ now gives
\[
8n+3=8(T_r+T_s+T_t)+3,\qquad
T_k=\frac{k(k+1)}2.
\]
Subtract $3$ and divide by $8$: $n=T_r+T_s+T_t$. Every nonnegative integer is therefore a sum of three triangular numbers.
\end{solution}

\begin{solution}{ex:zagier}
\idea The trick is a parity argument: an involution on a finite set has a number of fixed points with the same parity as the set. The set $S$ is designed so that one involution, $(x,y,z)\mapsto(x,z,y)$, has as fixed points exactly the decompositions $p=x^2+(2y)^2$; Zagier's second involution $f$ shows $|S|$ is odd. So the work is checking $f$ (routine algebra) and finding its fixed points.
\writeup (a) The equation $x^2+4yz=p$ bounds $x^2<p$ and $yz<p/4$. Since $x,y,z$ are positive integers, only finitely many triples are possible. The two boundaries in the definition of $f$ do not occur: $x=y-z$ would give $p=(y+z)^2$, and $x=2y$ would give $p=4y(y+z)$. Both contradict the primality of $p$. Thus every triple belongs to exactly one of the prescribed cases.

(b) In each case, the stated inequalities make the new coordinates positive. Substitution verifies that $x'^2+4y'z'=x^2+4yz=p$. For example, in the first case the image is $(x+2z,z,y-x-z)$, and
\[
(x+2z)^2+4z(y-x-z)=x^2+4yz.
\]
The other two cases expand in exactly the same way. The first case moves to the third, the third moves to the first, and the middle case stays in the middle; substituting twice returns the original triple. Hence $f$ is an involution.

(c) A fixed point cannot occur in the first or third case, since those exchange cases. In the middle case, the equality $f(x,y,z)=(x,y,z)$ forces $2y-x=x$, or $x=y$. Substituting into the defining equation gives
\[
p=x^2+4xz=x(x+4z).
\]
The primality of $p$ and positivity of $z$ force $x=y=1$, hence $z=(p-1)/4$. Thus there is exactly one fixed point. Every other element of the finite set $S$ belongs to a pair $\{s,f(s)\}$, so $|S|$ is odd.

(d) Swap the last two coordinates: $g(x,y,z)=(x,z,y)$. This is another involution of $S$. A finite set of odd size cannot be partitioned entirely into two-element orbits, so $g$ has a fixed point. There $y=z$, and the defining equation reads
\[
p=x^2+4y^2=x^2+(2y)^2.
\]
The oddness argument therefore produces the required two-square representation.
\end{solution}

\begin{solution}{ex:t1-which20}
\idea Theorem~\ref{thm:twosqn} is a condition on prime factors, not a search.
\writeup Check each candidate by the prime-factor criterion, and give a representation when it passes.

For $15=3\cdot5$, the prime $3\equiv3\pmod4$ has odd exponent, so $15$ is excluded. The number $18=2\cdot3^2$ passes, and $18=3^2+3^2$. The number $45=3^2\cdot5$ passes, and $45=3^2+6^2$. Finally $50=2\cdot5^2$ passes, with two immediate representations:
\[
50=5^2+5^2=7^2+1^2.
\]
The criterion decides existence; the displayed equalities verify the actual solutions.
\end{solution}

\section*{Sources and further reading}

\begin{itemize}
\item E.~Grosswald, \emph{Representations of Integers as Sums of Squares}, Springer, 1985. Everything about two, three, four and more squares, with $r_k(n)$ formulas.
\item D.~Zagier, A one-sentence proof that every prime $p\equiv1\pmod4$ is a sum of two squares, \emph{American Mathematical Monthly} 97 (1990), 144. The proof of Exercise~\ref{ex:zagier}.
\item D.~A.~Cox, \emph{Primes of the Form $x^2+ny^2$}, 3rd ed., AMS Chelsea, 2022, \S\S1--3. Fermat's questions about $x^2+2y^2$ and $x^2+3y^2$ (Exercises~\ref{ex:x2y2} and~\ref{ex:x3y2}) and how they lead to class field theory.
\item M.~N.~Huxley, Exponential sums and lattice points~III, \emph{Proceedings of the London Mathematical Society} 87 (2003), 591--609. The bound $O(R^{131/208+\varepsilon})$, about $R^{0.63}$, for the circle problem.
\item J.~Bourgain and N.~Watt, Mean square of zeta function, circle problem and divisor problem revisited, preprint, arXiv:1709.04340, 2017. A more recent improvement, to $R^{517/824+\varepsilon}\approx R^{0.6275}$, and a glimpse of current research.
\item A.~Weil, \emph{Number Theory: An Approach through History}, Birkh\"auser, 1984, Chapters~II--III. Fermat's letter to Mersenne of 25 December 1640, Euler's proof (letter to Goldbach, 1749) and Lagrange's four-square theorem (1770).
\end{itemize}
